\section{Conclusion}
\label{sec:conclusion}

We presented a complete pipeline for content- and preference-aware music recommendation grounded in a structured RDF knowledge graph. Two complementary KG variants encoding 7,013 tracks, 3,668 artists, and 285,037 users were validated through SPARQL inspection in GraphDB, confirming ontological alignment with Wikidata and structural soundness throughout.

A symbolic-feature autoencoder compressed 1,525-dimensional jSymbolic/music21 descriptors to a 128-dimensional bottleneck (final MSE: 0.616), and RotatE/ComplEx KG embeddings provided structural node initialisation across all entity types. Evaluation against three baselines established that KNN collaborative filtering achieves the strongest Overall Score (0.421) due to its near-complete catalogue coverage (91.6\%), while the popularity baseline attains slightly higher per-user NDCG at the cost of extreme popularity concentration. The XGBoost hybrid reveals the retrieval bottleneck when content features are used without graph structure.

The HGT recommender, trained with an intensity-weighted listwise loss with logit adjustment for popularity debiasing, achieves Recall@10~0.111 and Coverage~39.9\% in a consistent evaluation. Its slightly underperforming the KNN-CF in the unified benchmark is attributable to an undirected-training/directed-evaluation mismatch rather than a fundamental capacity limitation; retraining with consistent directed message passing is the straightforward remediation. Notably, the RotatE KGE initialisation does not improve over random initialisation, suggesting that KGE-learned structural features are subsumed by the HGT's own relational learning.

Future work should retrain the HGT with directed message passing, extend baselines to graph-collaborative methods such as LightGCN~\cite{he2020lightgcn} or KGCN~\cite{wang2019kgcn}, and conduct a systematic ablation isolating the contributions of KG structure, symbolic audio features, and the popularity-debiased loss.
